\section{Rangkuman}
\label{sec:rangkuman-bab11}

Pada bab ini telah dibahas:

\begin{itemize}
    \item Operasi logika bekerja pada setiap bit secara mandiri,
          tanpa menjalarkan \eng{carry} antarbit.

    \item \asm{AND} menghasilkan 1 bila kedua bit bernilai 1,
          \asm{OR} bila salah satu bit bernilai 1, \asm{XOR} bila
          kedua bit berbeda, dan \asm{NOT} membalik seluruh bit.

    \item \asm{AND}, \asm{OR}, dan \asm{XOR} selalu mengosongkan
          \flag{CF} dan \flag{OF}. \flag{ZF}, \flag{SF}, dan
          \flag{PF} mengikuti hasil. \asm{NOT} tidak mengubah
          \eng{flag}.

    \item \eng{Mask} dipakai untuk menyalakan bit dengan \asm{OR},
          memadamkan bit dengan \asm{AND} dan mask terbalik,
          membalik bit dengan \asm{XOR}, serta menguji bit dengan
          \asm{TEST}.

    \item \asm{TEST} melakukan \asm{AND} tanpa menyimpan hasil,
          sedangkan \asm{CMP} melakukan \asm{SUB} tanpa menyimpan
          hasil. \asm{TEST} dipakai untuk bit, \asm{CMP} untuk nilai
          keseluruhan.

    \item \asm{SHL} menggeser bit ke kiri dan mengisi bit kanan
          dengan 0. \asm{SHR} menggeser bit ke kanan dan mengisi bit
          kiri dengan 0. Bit yang keluar ditampung \flag{CF}.

    \item Pergeseran ke kiri setara dengan perkalian pangkat dua,
          sedangkan pergeseran ke kanan setara dengan pembagian
          pangkat dua untuk bilangan tidak bertanda.

    \item Penerapan operasi logika dan pergeseran meliputi
          pemeriksaan paritas, manipulasi bit, dan \eng{XOR cipher}
          sederhana.
\end{itemize}

\begin{checklistbox}
\vspace{1.1em}
Tandai setiap butir yang sudah Anda kuasai.

\begin{itemize}
    \item[\checkbox] Saya dapat menjelaskan tabel kebenaran
          \asm{AND}, \asm{OR}, \asm{XOR}, dan \asm{NOT}.
    \item[\checkbox] Saya dapat menjelaskan pengaruh operasi logika
          terhadap \eng{flag}.
    \item[\checkbox] Saya dapat menyalakan, memadamkan, membalik,
          dan menguji bit dengan mask.
    \item[\checkbox] Saya dapat membedakan \asm{TEST} dan
          \asm{CMP}.
    \item[\checkbox] Saya dapat menggunakan \asm{SHL} dan
          \asm{SHR}.
    \item[\checkbox] Saya dapat memakai pergeseran untuk perkalian
          dan pembagian pangkat dua.
    \item[\checkbox] Saya dapat menjelaskan program paritas dan
          \eng{XOR cipher}.
\end{itemize}

\vspace{0.4em}
Bila masih ada butir yang belum bertanda, ulangi bagian yang
bersesuaian sebelum melanjutkan ke bab berikutnya.
\end{checklistbox}

\begin{infobox}
\vspace{1.1em}
\textbf{Rujukan bab.} Materi pada bab ini diringkas dari
\cite{gunawan2019}, Bagian 7, hlm.~68--73, dan \cite{sto2001},
BAB X, hlm.~47--51.
\end{infobox}
